\scanpage{2}
\mspage{1}
\textcircled{1} Soit $\overline{\mathbf Q}$ la clôture algébrique de $\mathbf Q$ dans $\mathbf C$.
Si $\unsure{X=}\overline X$ est un schéma sur $\overline{\mathbf Q}$, on considère le
groupe $\mathcal G_{\overline X}$ des automorphismes de $\overline X$ en tant
que schéma absolu. On a une suite
\[
(1)\qquad 1\longrightarrow G_{\overline X}\longrightarrow\mathcal G_{\overline X}\longrightarrow\Gamma
\]
\margin{Question : \unsure{comme} \illegible{1} ?}
où
\[
(2)\qquad\left\{\begin{aligned}
\Gamma&=\operatorname{Gal}(\overline{\mathbf Q}/\mathbf Q),\\
G_{\overline X}&=\operatorname{Aut}(\overline X/\overline{\mathbf Q})
\end{aligned}\right.
\]
\margin{\diagram{\begin{tikzcd}[ampersand replacement=\&]
\overline{\mathbf Q}\arrow[r,hook]\arrow[d,phantom,"\cup" description]\&\Gamma(\unsure{X},\mathcal O_{\unsure{X}})\\
\mathbf Q\&
\end{tikzcd}}
\unsure{Il faut supposer} p.ex. que $\overline{\mathbf Q}$ est la
\unsure{clôt. int.} de $\mathbf Q$ ds $\unsure{\Gamma_X}$.}
Les sous-extensions $K$ de $\overline{\mathbf Q}/\mathbf Q$ correspondent
aux sous-groupes fermés $\Gamma_K$ de $\Gamma$ (d'indice
fini ssi $K$ fini sur $\mathbf Q$). À une restriction
du corps de base $\overline{\mathbf Q}$ de $\overline X$ à $K$ correspond
un relèvement
\[
(\unsure{4})\qquad
\diagram{\begin{tikzcd}[ampersand replacement=\&]
\mathcal G_{\overline X}\arrow[r,dash]\&\Gamma\arrow[d,phantom,"\cup" description]\\
\&\Gamma_K\arrow[ul]
\end{tikzcd}},
\]
\margin{\unsure{grain de sel d'effectivité de descente} \unsure{de gr. prof.}}
qui le définit si $\overline X$ est réduit\margin{?} et loc. de t.f.
sur $\overline{\mathbf Q}$. Cela montre que si $\overline X$ est de t.f.
\deleted{\unsure{quasi-proj.}}
sur $\overline{\mathbf Q}$, l'hom $\mathcal G_{\overline X}\mathbin{\unsure{-|\to}}\Gamma$ a une image
d'indice fini (puisqu'on peut \inserted{alors} réduire le
corps de base à une ext. finie $K$ de $\mathbf Q$).

\textcircled{2} \uline{Supposons $X$ connexe} \deleted{\unsure{et}}\inserted{$\ne\varnothing$}.
Considérons un revêtement universel $\widetilde X\xrightarrow{h}X$
\deleted{$X$}, et \deleted{considérons} posons
\[
(5)\qquad \mathcal E(X,\widetilde X)\ \unsure{\text{ou}}\ \mathcal E_X=\operatorname{Aut}(\widetilde X/X)
\]
\mspage{1}
